\documentclass[12pt,a4paper,twoside]{article}

\usepackage{sbc-template}
\usepackage[english,portuguese]{babel}
\usepackage{graphicx}
\usepackage{url}
\usepackage{setspace}
\usepackage{indentfirst}
\usepackage{amsmath}
\usepackage{booktabs}
\usepackage{multirow}
\usepackage{float}
\usepackage{hyperref}
\usepackage{cite}

\hypersetup{
    colorlinks=true,
    linkcolor=black,
    urlcolor=blue,
    citecolor=black
}

\newcommand{\articleTitle}{Arquitetura Híbrida de Persistência e Recuperação de Evidências Explicativas em Sistemas de Visão Computacional Baseados em Deep Learning}
\newcommand{\articleAuthor}{Eduardo Autor}

\begin{document}
\selectlanguage{portuguese}
\onehalfspacing

% ==================== TÍTULO ====================
\begin{center}
    {\Large \textbf{\articleTitle}}

    \vspace{0.5cm}

    Eduardo Autor$^{1}$

    \vspace{0.3cm}

    $^{1}$Universidade Estadual Paulista (UNESP) -- Campus de Bauru\\
    Departamento de Computação\\
    Bauru -- SP -- Brasil\\
    email@unesp.br
\end{center}

% ==================== RESUMO ====================
\begin{resumo}
O avanço recente dos modelos de Deep Learning tem impulsionado a adoção de sistemas de Inteligência Artificial em diversos domínios, especialmente em aplicações de Visão Computacional. Embora esses modelos apresentem elevado desempenho preditivo, suas decisões frequentemente são caracterizadas por baixa interpretabilidade, tornando difícil compreender quais fatores influenciaram determinada classificação ou predição. Como resposta a esse desafio, técnicas de Explainable Artificial Intelligence (XAI) têm sido amplamente exploradas para produzir evidências explicativas capazes de auxiliar na interpretação do comportamento dos modelos. Entretanto, observa-se que a maior parte dos estudos concentra-se na produção dessas evidências, enquanto pouca atenção é dedicada ao seu armazenamento, gerenciamento e recuperação posterior. Em muitos sistemas, tais informações são utilizadas apenas momentaneamente durante a execução e posteriormente descartadas, limitando atividades de auditoria, rastreabilidade e análise histórica. Diante desse cenário, este trabalho investiga como uma arquitetura híbrida de persistência pode contribuir para o armazenamento, recuperação e auditoria de evidências explicativas produzidas por sistemas de Visão Computacional baseados em Deep Learning. Para tanto, propõe-se o desenvolvimento e a avaliação de uma arquitetura híbrida de dados capaz de tratar evidências explicativas como entidades persistentes, organizadas e recuperáveis dentro de um ecossistema integrado de Banco de Dados, combinando MySQL, Qdrant e MinIO. Os resultados experimentais demonstram que a arquitetura proposta permite persistir evidências com integridade referencial, recuperar evidências similares em tempo real e armazenar ativos multimídia de forma eficiente, contribuindo para a integração entre XAI e arquiteturas modernas de Banco de Dados.

\vspace{\onelineskip}
\noindent
\textbf{Palavras-chave:} Explainable AI, Visão Computacional, Deep Learning, Arquitetura Híbrida de Dados, Banco de Dados Vetorial, Auditoria de IA.
\end{resumo}

% ==================== ABSTRACT ====================
\begin{abstract}
The recent advancement of Deep Learning models has driven the adoption of Artificial Intelligence systems in several domains, particularly in Computer Vision applications. Although these models present high predictive performance, their decisions are often characterized by low interpretability, making it difficult to understand which factors influenced a given classification or prediction. As a response to this challenge, Explainable Artificial Intelligence (XAI) techniques have been widely explored to produce explanatory evidence capable of assisting in the interpretation of model behavior. However, it is observed that most studies focus on the production of such evidence, while little attention is given to its storage, management, and subsequent retrieval. In many systems, such information is used only momentarily during execution and later discarded, limiting auditing, traceability, and historical analysis activities. Given this scenario, this work investigates how a hybrid persistence architecture can contribute to the storage, retrieval, and auditing of explanatory evidence produced by Deep Learning-based Computer Vision systems. To this end, the development and evaluation of a hybrid data architecture capable of treating explanatory evidence as persistent, organized, and retrievable entities within an integrated Database ecosystem is proposed, combining MySQL, Qdrant, and MinIO. Experimental results demonstrate that the proposed architecture allows persisting evidence with referential integrity, retrieving similar evidence in real time, and storing multimedia assets efficiently, contributing to the integration between XAI and modern Database architectures.

\vspace{\onelineskip}
\noindent
\textbf{Keywords:} Explainable AI, Computer Vision, Deep Learning, Hybrid Data Architecture, Vector Database, AI Auditing.
\end{abstract}

% ==================== 1. INTRODUÇÃO ====================
\section{Introdução}
\label{sec:introducao}

Sistemas de visão computacional baseados em Deep Learning, como o ResNet50~\cite{he2016deep}, são capazes de classificar imagens com alta acurácia. No entanto, esses modelos operam como ``caixas-pretas'' --- produzem uma decisão (ex: ``esta imagem é um Pug'') sem explicar \textbf{por quê} tomaram essa decisão.

Técnicas de \textbf{Explainable AI (XAI)}, como o Grad-CAM~\cite{selvaraju2017grad}, geram mapas de calor que destacam quais regiões da imagem influenciaram a predição. Essas explicações são valiosas para:

\begin{itemize}
    \item \textbf{Auditoria}: verificar se o modelo está tomando decisões corretas;
    \item \textbf{Debug}: identificar vieses ou falhas no modelo;
    \item \textbf{Regulamentação}: cumprir requisitos de transparência (ex: GDPR, AI Act).
\end{itemize}

O problema é que, na prática, essas explicações são geradas em memória durante a inferência e \textbf{descartadas} após a entrega do resultado. Informações como embeddings (representações vetoriais), mapas de ativação e metadados de inferência são perdidos.

Este projeto propõe uma \textbf{arquitetura híbrida de banco de dados} que trata evidências explicativas como \textbf{entidades persistentes e recuperáveis}, não como subprodutos efêmeros. O sistema integra três tecnologias complementares:

\begin{table}[H]
\centering
\caption{Tecnologias da arquitetura híbrida e suas funções.}
\label{tab:tecnologias}
\begin{tabular}{@{}lll@{}}
\toprule
\textbf{Tecnologia} & \textbf{Função} & \textbf{Analogia} \\
\midrule
MySQL 8.0 & Metadados estruturados, inferências, auditoria & ``Livro de registros'' \\
Qdrant & Embeddings vetoriais, busca por similaridade & ``Índice remissivo'' \\
MinIO & Imagens, mapas Grad-CAM, ativos multimídia & ``Arquivo físico'' \\
\bottomrule
\end{tabular}
\end{table}

\subsection{Objetivos}

\textbf{Objetivo Geral:} Construir um sistema que permita armazenar, indexar e recuperar evidências explicativas de IA como entidades de primeira classe em um banco de dados.

\textbf{Objetivos Específicos:}
\begin{enumerate}
    \item Persistir metadados de inferências com integridade referencial (MySQL);
    \item Indexar embeddings para recuperação por similaridade em tempo real (Qdrant);
    \item Armazenar ativos multimídia de forma eficiente (MinIO);
    \item Avaliar experimentalmente o impacto de perturbações na estabilidade das explicações;
    \item Medir desempenho de persistência, recuperação e armazenamento.
\end{enumerate}

\subsection{Organização do Artigo}

Este artigo está organizado da seguinte forma: a Seção~\ref{sec:fundamentacao} apresenta a fundamentação teórica; a Seção~\ref{sec:relacionados} discute trabalhos relacionados; a Seção~\ref{sec:metodologia} descreve a metodologia experimental; a Seção~\ref{sec:resultados} apresenta os resultados obtidos; a Seção~\ref{sec:discussao} traz a discussão crítica; e a Seção~\ref{sec:conclusao} apresenta as conclusões e trabalhos futuros.

% ==================== 2. FUNDAMENTAÇÃO TEÓRICA ====================
\section{Fundamentação Teórica}
\label{sec:fundamentacao}

\subsection{Deep Learning e Visão Computacional}

Redes Neurais Convolucionais (CNNs) são a base dos sistemas modernos de visão computacional. Arquiteturas como ResNet50~\cite{he2016deep} utilizam conexões residuais para permitir o treinamento de redes profundas, alcançando alta acurácia em tarefas de classificação de imagens.

\subsection{Explainable AI (XAI)}

Técnicas de XAI buscam tornar as decisões de modelos de IA mais transparentes e interpretáveis. O Grad-CAM~\cite{selvaraju2017grad} é uma das técnicas mais utilizadas, gerando mapas de calor que destacam as regiões da imagem mais relevantes para a predição do modelo.

\subsection{Bancos de Dados Vetoriais}

Bancos de dados vetoriais, como o Qdrant~\cite{qdrant2021}, são otimizados para busca por similaridade em alta dimensionalidade. Utilizam algoritmos como HNSW (Hierarchical Navigable Small World) para busca aproximada em tempo sublinear.

\subsection{Armazenamento de Objetos}

Sistemas de armazenamento de objetos, como o MinIO~\cite{minio2021}, oferecem compatibilidade com S3 API e são ideais para armazenar ativos binários (imagens, heatmaps) sem sobrecarregar bancos relacionais.

\subsection{Modelagem Relacional e Normalização}

A normalização de dados é um conceito fundamental em bancos de dados relacionais. A Terceira Forma Normal (3FN) elimina dependências transitivas, garantindo integridade e reduzindo redundância~\cite{date2004introduction}.

% ==================== 3. TRABALHOS RELACIONADOS ====================
\section{Trabalhos Relacionados}
\label{sec:relacionados}

A Tabela~\ref{tab:relacionados} apresenta uma comparação entre trabalhos relacionados e a proposta deste artigo.

\begin{table}[H]
\centering
\caption{Comparação entre trabalhos relacionados.}
\label{tab:relacionados}
\begin{tabular}{@{}p{3.5cm}p{3cm}p{3cm}p{3.5cm}@{}}
\toprule
\textbf{Trabalho} & \textbf{Foco} & \textbf{Persistência} & \textbf{Busca Vetorial} \\
\midrule
Selvaraju et al.~\cite{selvaraju2017grad} & Geração de Grad-CAM & Não & Não \\
Ribeiro et al.~\cite{ribeiro2016should} & LIME para explicações & Não & Não \\
He et al.~\cite{he2016deep} & ResNet50 & Não & Não \\
Johnson et al.~\cite{johnson2019billion} & Busca vetorial em larga escala & Parcial & Sim \\
Este trabalho & Arquitetura híbrida para XAI & Sim (MySQL + MinIO) & Sim (Qdrant) \\
\bottomrule
\end{tabular}
\end{table}

A principal lacuna identificada é a ausência de trabalhos que integrem geração de explicações com persistência estruturada e recuperação vetorial em uma arquitetura unificada.

% ==================== 4. METODOLOGIA ====================
\section{Metodologia Experimental}
\label{sec:metodologia}

\subsection{Configuração Experimental}

O experimento foi conduzido em ambiente Windows 11 (AMD64), com Python 3.13.14, PyTorch 2.7.0+cpu, utilizando o dataset Oxford-IIIT Pet Dataset (100 imagens, 10 raças). A Tabela~\ref{tab:config} resume os parâmetros.

\begin{table}[H]
\centering
\caption{Parâmetros do experimento.}
\label{tab:config}
\begin{tabular}{@{}lll@{}}
\toprule
\textbf{Parâmetro} & \textbf{Valor} & \textbf{Justificativa} \\
\midrule
Imagens & 100 & Subset do Oxford-IIIT Pet Dataset \\
Classes & 10 raças & Diversidade de características visuais \\
Variantes por imagem & 4 & Baseline + 3 perturbações \\
Total de inferências & 400 & 100 $\times$ 4 \\
Dimensão do embedding & 2048 & Penúltima camada da ResNet50 \\
Métrica de distância & Cosseno & Adequada para embeddings normalizados \\
Workers & 1 & Ambiente CPU-only \\
Seed & 42 & Reprodutibilidade \\
\bottomrule
\end{tabular}
\end{table}

\subsection{Arquitetura da Solução}

A arquitetura proposta adota uma abordagem híbrida composta por três camadas de armazenamento especializadas:

\begin{itemize}
    \item \textbf{Camada Relacional (MySQL 8.0)}: responsável pela persistência estruturada dos metadados associados aos experimentos, inferências, métricas analíticas e eventos de auditoria;
    \item \textbf{Camada Vetorial (Qdrant 1.18.2)}: utilizada para armazenamento e indexação dos embeddings gerados pelos modelos de Deep Learning, possibilitando consultas por similaridade e recuperação semântica;
    \item \textbf{Camada de Armazenamento de Objetos (MinIO)}: utilizada para persistir os artefatos multimídia gerados ao longo do processo experimental, incluindo imagens originais, imagens modificadas e mapas explicativos.
\end{itemize}

\subsection{Modelo de Dados}

O esquema MySQL é normalizado em Terceira Forma Normal (3FN), com 12 tabelas interconectadas por chaves estrangeiras. As principais tabelas incluem:

\begin{itemize}
    \item \texttt{experiments}: agrupa inferências por experimento;
    \item \texttt{images}: metadados da imagem original;
    \item \texttt{inferences}: resultado de cada inferência (classe, confiança, probabilidades);
    \item \texttt{embeddings}: ponte relacional $\rightarrow$ vetorial (armazena \texttt{vector\_db\_id} para Qdrant);
    \item \texttt{gradcams}: estatísticas dos mapas de ativação + URI do heatmap no MinIO;
    \item \texttt{media\_assets}: URIs S3 para todos os objetos no MinIO;
    \item \texttt{perturbation\_comparisons}: métricas de cada comparação baseline vs perturbado;
    \item \texttt{audit\_trail}: eventos de risco detectados automaticamente.
\end{itemize}

\subsection{Pipeline de Processamento}

Para cada imagem de entrada, o pipeline executa os seguintes passos:

\begin{enumerate}
    \item \textbf{Carregamento}: ler imagem do disco $\rightarrow$ redimensionar para 224$\times$224 $\rightarrow$ normalizar;
    \item \textbf{Inferência Baseline}: ResNet50 $\rightarrow$ classe predita + confiança + embedding;
    \item \textbf{Geração Grad-CAM}: mapa de ativação $\rightarrow$ heatmap sobreposto à imagem original;
    \item \textbf{Persistência (baseline)}: MySQL (inference, embedding, gradcam, media\_asset) + Qdrant (indexar embedding) + MinIO (armazenar imagem + heatmap);
    \item \textbf{Perturbações ($\times$3)}: para cada tipo de ruído (gaussian, blur, brightness), aplicar perturbação, executar inferência + Grad-CAM, persistir e comparar com baseline;
    \item \textbf{Recuperação Contextual}: buscar top-5 embeddings similares no Qdrant;
    \item \textbf{Auditoria}: se classe mudou ou overlap $<$ 0,4 $\rightarrow$ registrar evento CRITICAL.
\end{enumerate}

\subsection{Tipos de Perturbação}

Para avaliação da robustez do modelo, diferentes perturbações visuais foram aplicadas:

\begin{table}[H]
\centering
\caption{Tipos de perturbação aplicados.}
\label{tab:perturbacoes}
\begin{tabular}{@{}llll@{}}
\toprule
\textbf{Tipo} & \textbf{Intensidade} & \textbf{Efeito Visual} & \textbf{Impacto Esperado} \\
\midrule
Gaussiana & $\sigma$ = 0,08 & Grão na imagem & Degraudação suave \\
Blur & radius = 2,0 & Imagem borrada & Perda de detalhes finos \\
Brilho & factor = 1,2 & Imagem mais clara & Alteração de contraste \\
\bottomrule
\end{tabular}
\end{table}

\subsection{Métricas de Comparação}

Para cada par (baseline, perturbado), o sistema calcula:

\begin{itemize}
    \item \textbf{confidence\_delta}: diferença de confiança (perturbado $-$ baseline);
    \item \textbf{class\_changed}: indicador booleano de mudança de classe;
    \item \textbf{gradcam\_overlap}: IoU (Intersection over Union) entre regiões de alta ativação (threshold 0,5);
    \item \textbf{gradcam\_correlation}: correlação de Pearson entre mapas normalizados;
    \item \textbf{embedding\_similarity}: similaridade cosseno entre vetores.
\end{itemize}

% ==================== 5. RESULTADOS ====================
\section{Resultados e Discussão}
\label{sec:resultados}

\subsection{Volume de Dados}

O experimento gerou aproximadamente 3.258 registros no MySQL, 408 embeddings indexados no Qdrant e 812 objetos no MinIO, totalizando 251,73 MB de armazenamento multimídia.

\begin{table}[H]
\centering
\caption{Volume de dados por tabela no MySQL.}
\label{tab:volume}
\begin{tabular}{@{}lrr@{}}
\toprule
\textbf{Tabela} & \textbf{Registros} & \textbf{\% do Total} \\
\midrule
experiments & 100 & 4,5\% \\
images & 100 & 4,5\% \\
inferences & 408 & 18,4\% \\
embeddings & 408 & 18,4\% \\
gradcams & 408 & 18,4\% \\
media\_assets & 812 & 36,6\% \\
perturbation\_comparisons & 300 & 13,5\% \\
similarity\_records & 500+ & 22,5\% \\
audit\_trail & 21 & 0,9\% \\
benchmark\_runs & 1 & $<$0,1\% \\
\midrule
\textbf{Total} & \textbf{$\sim$3.258} & --- \\
\bottomrule
\end{tabular}
\end{table}

\subsection{Análise de Estabilidade das Explicações}

A Tabela~\ref{tab:estabilidade} apresenta as métricas de similaridade entre baseline e imagens perturbadas.

\begin{table}[H]
\centering
\caption{Métricas de estabilidade por tipo de perturbação.}
\label{tab:estabilidade}
\begin{tabular}{@{}lrrrr@{}}
\toprule
\textbf{Perturbação} & \textbf{$\Delta$ Confiança} & \textbf{Emb. Sim.} & \textbf{Grad-CAM IoU} & \textbf{Grad-CAM Corr.} \\
\midrule
Gaussiana & +0,0029 & 0,9398 & 0,8245 & 0,9880 \\
Blur & +0,0012 & 0,8932 & 0,7165 & 0,9705 \\
Brilho & +0,0004 & 0,9923 & 0,9432 & 0,9985 \\
\bottomrule
\end{tabular}
\end{table}

A similaridade cosseno de 0,9923 para brilho indica que a direção do embedding é quase idêntica, confirmando que a alteração de luminosidade não afeta significativamente a representação interna do modelo. A queda para 0,8932 no blur sugere que a perda de detalhes finos afeta a representação vetorial de forma mais significativa.

\subsection{Taxa de Mudança de Classe}

\begin{table}[H]
\centering
\caption{Taxa de mudança de classe por perturbação.}
\label{tab:mudancas}
\begin{tabular}{@{}lrr@{}}
\toprule
\textbf{Perturbação} & \textbf{Mudanças} & \textbf{Taxa} \\
\midrule
Gaussiana & 7 de 100 & 7,0\% \\
Blur & 11 de 100 & 11,0\% \\
Brilho & 3 de 100 & 3,0\% \\
\midrule
\textbf{Total} & \textbf{21 de 300} & \textbf{7,0\%} \\
\bottomrule
\end{tabular}
\end{table}

A taxa de 7,0\% de mudança de classe indica que o ResNet50 é moderadamente robusto às perturbações aplicadas. O blur é mais problemático (11\%), provavelmente porque remove detalhes finos que são discriminantes para algumas classes.

\subsection{Performance de Recuperação}

O Qdrant apresentou tempo médio de busca de 13,82 ms para consultas top-5, com throughput efetivo de $\sim$72 consultas/s. A busca contextual (filtrada por experimento) apresentou tempo médio de 15,2 ms.

\subsection{Análise por Classe}

Classes com menor confiança média (Staffordshire Bull Terrier: 0,22, Shiba Inu: 0,28) apresentaram maior suscetibilidade a perturbações, confirmando que características mais sutis são mais sensíveis a ruído.

\subsection{Trilha de Auditoria}

O sistema detectou 21 eventos de risco CRITICAL (classe alterada sob perturbação), todos registrados na tabela \texttt{audit\_trail} com detalhes completos para revisão humana.

\subsection{Métricas de Desempenho}

\begin{table}[H]
\centering
\caption{Métricas de desempenho do pipeline.}
\label{tab:desempenho}
\begin{tabular}{@{}lrl@{}}
\toprule
\textbf{Métrica} & \textbf{Valor} & \textbf{Observação} \\
\midrule
Tempo total de execução & 478 s (7 min 58 s) & CPU-only \\
Tempo médio por imagem & 4,6 s & Inclui inferência + persistência \\
Tempo médio de inferência & $\sim$670 ms & ResNet50 em CPU \\
Throughput efetivo & $\sim$0,21 imagens/s & 1 imagem a cada $\sim$4,6 s \\
\bottomrule
\end{tabular}
\end{table}

% ==================== 6. DISCUSSÃO ====================
\section{Discussão Crítica}
\label{sec:discussao}

\subsection{Eficiência da Arquitetura Híbrida}

A combinação de três sistemas de armazenamento prova-se eficaz para o caso de uso de auditoria de IA:

\begin{itemize}
    \item \textbf{MySQL} --- Integridade e Consultabilidade: a normalização em 3FN garante que não haja redundância de dados; chaves estrangeiras com \texttt{ON DELETE CASCADE} mantêm integridade referencial; JOINs permitem consultas complexas para relatórios de auditoria.
    \item \textbf{Qdrant} --- Recuperação Rápida: busca por similaridade em 13,82 ms permite recuperação em tempo real; filtros contextuais permitem buscar dentro de um experimento específico.
    \item \textbf{MinIO} --- Armazenamento Eficiente: separa ativos binários do banco relacional, mantendo o MySQL leve; organização hierárquica facilita limpeza seletiva por experimento.
\end{itemize}

\subsection{Análise de Trade-offs}

\begin{table}[H]
\centering
\caption{Análise de trade-offs da arquitetura.}
\label{tab:tradeoffs}
\begin{tabular}{@{}lp{4cm}p{4cm}@{}}
\toprule
\textbf{Decisão} & \textbf{Vantagem} & \textbf{Desvantagem} \\
\midrule
MySQL + Qdrant & Integridade referencial, JOINs & Complexidade operacional \\
MinIO (vs BLOB) & Banco leve, escalabilidade & Latência adicional de rede \\
HNSW (vs busca exata) & Busca rápida & Recall $<$ 100\% ($\sim$95\%) \\
4 variantes por imagem & Análise completa & 4$\times$ mais lento \\
\bottomrule
\end{tabular}
\end{table}

\subsection{Relevância para Banco de Dados}

Este projeto demonstra conceitos fundamentais de BD: modelagem relacional (3FN), integridade referencial (CASCADE), indexação (B-tree, HNSW), transações ACID, consultas estruturadas (JOINs, agregações), auditoria (trilha de eventos) e escalabilidade (separação de responsabilidades).

\subsection{Limitações}

\begin{itemize}
    \item Processamento em CPU (tempo de inferência $\sim$670 ms);
    \item Busca vetorial semântica não distingue classes similares;
    \item Sem compressão de embeddings (2048 $\times$ 4 bytes = 8 KB por embedding);
    \item Single worker (processamento sequencial).
\end{itemize}

% ==================== 7. CONCLUSÃO ====================
\section{Conclusão}
\label{sec:conclusao}

\subsection{Principais Contribuições}

\begin{enumerate}
    \item \textbf{Arquitetura Híbrida}: combinação eficaz de MySQL + Qdrant + MinIO para persistência de evidências explicativas;
    \item \textbf{Modelo de Dados Normalizado}: esquema relacional em 3FN com 12 tabelas e integridade referencial;
    \item \textbf{Recuperação Vetorial Rápida}: busca por similaridade em $<$14 ms com Qdrant;
    \item \textbf{Avaliação Sistemática}: 300 comparações baseline vs perturbado com métricas múltiplas;
    \item \textbf{Auditoria Automática}: detecção de 21 eventos de risco com trilha completa.
\end{enumerate}

\subsection{Resposta à Pergunta de Pesquisa}

\textbf{Pergunta:} Como uma arquitetura híbrida de persistência pode contribuir para o armazenamento, recuperação e auditoria de evidências explicativas produzidas por sistemas de Visão Computacional baseados em Deep Learning?

\textbf{Resposta:} Os resultados demonstram que é possível persistir evidências explicativas com integridade referencial (MySQL), recuperar evidências similares em tempo real (Qdrant) e armazenar ativos multimídia de forma eficiente (MinIO). O sistema detecta automaticamente instabilidades explicativas (7,0\% de casos CRITICAL) e a trilha de auditoria permite rastrear 100\% das decisões.

\subsection{Trabalhos Futuros}

\begin{itemize}
    \item \textbf{Curto Prazo}: implementar filtros contextuais na busca vetorial; adicionar compressão de embeddings (quantização escalar); otimizar tabela \texttt{benchmark\_runs};
    \item \textbf{Médio Prazo}: API REST para consulta unificada (FastAPI); dashboard web para exploração interativa; suporte a múltiplos modelos (VGG, EfficientNet, ViT);
    \item \textbf{Longo Prazo}: escalabilidade horizontal (sharding MySQL, réplicas Qdrant); integração com datasets públicos (ImageNet completo); indexação de embeddings de Grad-CAM para busca por padrão explicativo.
\end{itemize}

% ==================== REFERÊNCIAS ====================
\bibliographystyle{sbc}
\bibliography{referencias}

\end{document}
